\subsection{Simple canonical-form weights}

\begin{theorem}[Copy independence for a non-nilpotent sector decomposition]
    \label{thm:mpdo_simple_weight_copy_independent}
    \lean{MPOTensor.weight_copy_independent_of_isSourceZCL}
    \leanok
    \uses{def:mpo_source_zcl, def:mpdo_doubled_phys_trace_transfer,
        lem:mpdo_doubled_phys_trace_transfer_to_mps,
        lem:mpdo_doubled_phys_trace_transfer_blockwise}
    Let a tensor with zero correlation length be gauge-equivalent, through
    block-diagonal copy gauges, to an assembled sector decomposition whose
    basis elements all have non-nilpotent physical-trace transfer. Then the
    weights do not depend on the copy index: $\mu_{j,q}=\mu_{j}$
    \cite[Appendix~C.2, lines~1646--1661]{Cirac2016MPDO_arXiv}.
\end{theorem}

\begin{proof}\leanok
    Zero correlation length states
    $\mathcal T_M^{2}=\lambda\,\mathcal T_M$ with $\lambda>0$.  If
    $\mathcal T_M=X\mathcal T_SX^{-1}$ is the block-diagonal gauge to the
    transfer of the assembled sector decomposition, then cancellation of
    $X$ and $X^{-1}$ gives
    \[
        \mathcal T_S^{2}=\lambda\,\mathcal T_S.
    \]
    By
    Lemmas~\ref{lem:mpdo_doubled_phys_trace_transfer_to_mps}
    and~\ref{lem:mpdo_doubled_phys_trace_transfer_blockwise} the identity
    descends to every copy:
    \[
        \mu_{j,q}^{2}\,\mathcal T_{A_{j}}^{2}
        =\lambda\,\mu_{j,q}\,\mathcal T_{A_{j}},
        \qquad\text{hence}\qquad
        \mathcal T_{A_{j}}^{2}
        =\frac{\lambda}{\mu_{j,q}}\,\mathcal T_{A_{j}}.
    \]
    The left side does not depend on $q$, and
    $\mathcal T_{A_{j}}\ne0$ because it is not nilpotent, so
    $\lambda/\mu_{j,q}$ is independent of $q$, which gives
    $\mu_{j,q}=\mu_{j,q'}$.
\end{proof}

\begin{corollary}[Simple canonical-form tensors with zero correlation length]
    \label{cor:mpdo_simple_zcl_weight_copy_independent}
    \lean{MPOTensor.IsSimpleCanonicalForm.exists_weight_copy_independent_of_isSourceZCL}
    \leanok
    \uses{def:mpdo_simple_cf_tensor, def:mpo_source_zcl,
        thm:mpdo_simple_weight_copy_independent}
    A simple canonical-form tensor
    (Definition~\ref{def:mpdo_simple_cf_tensor}) with zero correlation length
    has non-nilpotent basis transfers and weights independent of the copy
    index.
\end{corollary}

\begin{proof}\leanok
    Apply Theorem~\ref{thm:mpdo_simple_weight_copy_independent} to the
    witness supplied by the simple canonical-form predicate.
\end{proof}

\begin{definition}[Common sector weight and absorbed representative]
    \label{def:mpdo_common_sector_weight_absorption}
    \lean{MPSTensor.SectorDecomposition.commonWeight}
    \lean{MPSTensor.SectorDecomposition.commonWeightAbsorbedBasis}
    \leanok
    \uses{def:sector_weight_data}
    Suppose the weights over each representative do not depend on the copy
    index. Choose the common weight $\mu_j$ from any copy over $j$ and absorb
    it into the representative:
    \[
        \mathcal K_j=\mu_j A_j.
    \]
    This is the convention of
    \cite[Appendix~C.2, equation~CFK, lines~1660--1665]
      {Cirac2016MPDO_arXiv}.
\end{definition}

\begin{lemma}[Every copy has the common weight]
    \label{lem:mpdo_copy_weight_eq_common_weight}
    \lean{MPSTensor.SectorDecomposition.weight_eq_commonWeight}
    \leanok
    \uses{def:mpdo_common_sector_weight_absorption}
    For every copy $q$ over $j$, one has $\mu_{j,q}=\mu_j$.
\end{lemma}

\begin{proof}\leanok
    Compare the $q$-th copy with the distinguished copy used to define
    $\mu_j$.
\end{proof}

\begin{theorem}[Coefficient after common-weight absorption]
    \label{thm:mpdo_common_weight_coefficient}
    \lean{MPSTensor.SectorDecomposition.coeff_eq_copies_mul_commonWeight_pow}
    \leanok
    \uses{def:mpdo_common_sector_weight_absorption,
      def:sector_weight_data, lem:mpdo_copy_weight_eq_common_weight}
    Let $r_j$ be the number of horizontal copies of representative $j$. Then
    its coefficient on a chain of length $N$ is
    \[
        \sum_q\mu_{j,q}^{N}=r_j\mu_j^{N}.
    \]
    Thus the natural multiplicity is the horizontal copy count $r_j$.
\end{theorem}

\begin{proof}\leanok
    Every term in the copy sum equals $\mu_j^N$, and there are $r_j$ terms.
\end{proof}

\begin{theorem}[Canonical form with natural multiplicities]
    \label{thm:mpdo_common_weight_absorbed_mpv}
    \lean{MPSTensor.SectorDecomposition.mpv_toTensor_eq_sum_copies_mul_commonWeightAbsorbedBasis}
    \leanok
    \uses{def:mpdo_common_sector_weight_absorption,
      lem:paper_decomp_formula, lem:mpv_smul,
      thm:mpdo_common_weight_coefficient}
    After absorbing the common weight, the matrix product vector of the
    assembled tensor is
    \[
        \mathcal V_N(S)=\sum_j r_j\,\mathcal V_N(\mathcal K_j).
    \]
    This is equation~CFK of
    \cite[Appendix~C.2, lines~1660--1665]{Cirac2016MPDO_arXiv}; the same
    multiplicities occur in equation~sigmaNKj of
    \cite[Appendix~C.2, lines~1756--1759]{Cirac2016MPDO_arXiv}.
\end{theorem}

\begin{proof}\leanok
    Expand the assembled tensor by representatives, use
    Theorem~\ref{thm:mpdo_common_weight_coefficient}, and absorb
    $\mu_j^N$ into the length-$N$ matrix product vector of $A_j$.
\end{proof}
